\documentclass[9pt,t]{beamer}

\usepackage[lang=fr, theme=ocots, math={analysis}, institution={n7}]{ocots}

\title{Suites réelles}
\author{ocourses}
\date{2026--2027}

\begin{document}

\slidechapter{1}{Suites réelles}
\slidetitlepage

%------------------------------------------------------------------
\begin{slide}{Convergence}

Une suite converge si ses termes finissent par rester aussi près qu'on le
veut d'un réel fixé.

\begin{definition}[title=Suite convergente]
    $(u_n)$ converge vers $\ell$ si
    $\forall \varepsilon > 0,\ \exists N,\ \forall n \ge N,\ \abs{u_n - \ell} \le \varepsilon$.
\end{definition}

\end{slide}

%------------------------------------------------------------------
\begin{slide}{Unicité}

\begin{theorem}[title=Unicité]
    La limite d'une suite convergente est unique.
\end{theorem}

\end{slide}

%------------------------------------------------------------------
\begin{slide}{Limite monotone}

Pour une suite croissante, la borne supérieure fournit la limite sans
qu'on ait à la deviner.

\begin{theorem}[title=Limite monotone]
    Toute suite croissante et majorée converge.
\end{theorem}

\end{slide}

\end{document}
